\section{A complete autoregressive decoder}
\label{sec:full-decoder-main}

This section states the decoder bound announced in the introduction
and the ideas of its proof. The index of Section~\ref{sec:index-main}
assumes finite records. The decoder also constructs each new key and
value with the neural model and charges that work.

Fix a dyadic embedding table with entries of magnitude at most $s$
and a causal network of the pre-norm form~\eqref{eq:main-prenorm}.
Use dyadic residual steps in $[0,1]$; convex residual mixtures with
coefficients in $[0,1]$ are also allowed. Allow an affine attention
output map and an affine vocabulary head. Parameters have at most $b$ bits,
row norms and biases are at most $s\ge1$, the temperature is
nonnegative dyadic or $\sqrt n$, and every normalization
has positive stabilizer at least $2^{-b}$. Each layer's fixed
linear key map has rank at most $r$. Its keys at all positions
must lie in one common affine space; separate rank bounds on
position-dependent maps do not imply that promise.

\begin{theorem}[Exact decoding with maintained neural keys]
\label{thm:main-decoder}
For each fixed $r$, there is a constructive polynomial $F_r$
such that a prompt of length $T_0\ge1$ can be prefilled in
\[
 T_0F_r(n,D,V,b,s,\beta,\log(T_0+2))
\]
bit operations. Every subsequent token has conditional expected
cost at most
\[
 F_r(n,D,V,b,s,\beta,\log(T+2)),
\]
given every previously sampled token history. This includes key
and value creation, insertion, random bits, and all refinements.
The output has the exact original autoregressive law.
At $\beta=\sqrt n$, with $D,V,b,s$ polynomially bounded in $n$,
the per-token bound is $\poly_r(n,\log T)$.
The affine head's logit range is determined from the parameter
bounds. The polynomial degree depends on $r$.
\end{theorem}

Theorem~\ref{thm:momentdecoder} proves this statement for logits
promised to lie in $[-1,1]$, and Corollary~\ref{cor:affineheadmoment}
removes that promise for the affine head.

In a bounded chart, write a score, up to a common offset, as
$\langle t,z_j\rangle$, where $z_j\in[-1,1]^r$ and
$\|t\|_1\le H=\poly_r(n,s,\beta)$. A degree-$m$ Taylor
expansion uses $\binom{m+r}{r}$ monomials. Store their sums
and their value-weighted sums as exact integers over a common
dyadic denominator. Each new record adds to these counters.
Degree $m=O(H+p)$ yields certified attention accuracy $2^{-p}$;
the average denominator is at least $e^{-H}$, so this includes
the error of division. At $\beta=\sqrt n$, rational even and
odd degree sums keep the specified square root exact.

Attention is Lipschitz from the maximum score error to the total
probability error with constant two, independently of the number
of keys. For positive RMS stabilizers, the logarithm of the Lipschitz
constant is of order $b+\log n$. Consequently errors propagate through
the $D$ layers, with no linear factor in the number of positions.
A precision of
$p+O\bigl((D+1)[b+\log(n+s+\beta+2)]+\log(D+V+2)\bigr)$
suffices.
Rounding the normalized stream before its exact key-map product
preserves the promised rank. This detail prevents a finite
approximation from silently changing the structural hypothesis.

Routine precision is a sufficiently large multiple of $\log T$.
If a uniform draw is too close to a categorical boundary, replay
the actual token history at higher precision until the sample is
certified. An unresolved precision-$p$ comparison has probability
at most $4V2^{-p}$. Choosing routine error of order $T^{-3}/V$
pays for a complete replay on that rare branch. A second routine
index replays two old tokens after each append and catches up
before the current precision expires. Ordinary maintenance uses
at most three deterministic token evaluations per step.
Figure~\ref{fig:decoder-overview} shows this maintenance loop.

For a basic tanh model with an affine head that is fixed in all
parameters, the key ranks are fixed even when the key matrices have
full rank. The explicit implementation gives
$O_{\rm model}(\log^{2r+12}(T+2))$ expected work per token, with
model-dependent constants that can be enormous
(Corollary~\ref{cor:fixedmodelcontext}).
Fixed coordinate-pair positional rotations also admit a
$\log^{O_{\rm model}(1)}T$ bound: certified trigonometric
evaluation costs polynomial time in precision and $\log T$
(Corollary~\ref{cor:fixedpositionalrotation}). The same
finite-evaluation argument covers specified GELU and SwiGLU updates
using their polynomial range and Lipschitz bounds on bounded
normalized inputs (Corollary~\ref{cor:momentgeluswiglu}).
